\subsubsection{Partícula en plano horizontal acoplada a resorte vertical}

Por un plano horizontal desliza sin roce una partícula de masa $m$ unida a un hilo. Éste pasa por un agujero y termina unido a un resorte de constante elástica $k$ ubicado verticalmente debajo del agujero. Cuando el resorte está en su largo natural, la partícula se encuentra justo en el agujero. Se estudia la dinámica de la partícula cuando es soltada a una distancia $r_0$ del agujero con una velocidad inicial perpendicular al hilo de magnitud $v_0$ en la dirección $\vec{e}_\theta$. Asuma que la gravedad actúa hacia abajo.

\begin{enumerate}
    \item Defina los marcos de referencia y sistemas de coordenadas adecuados (se recomiendan coordenadas polares en el plano).
    \item Analice la cinemática de la partícula y determine las expresiones para la velocidad y aceleración.
    \item Analice la cinética del sistema para obtener las ecuaciones diferenciales de movimiento utilizando la Segunda Ley de Newton o el formalismo de Lagrange.
\end{enumerate}

\begin{figure}[!ht]
    \centering
    % --- VISTA A: PERSPECTIVA GENERAL ---
    \begin{minipage}{0.48\textwidth}
        \centering
        \begin{tikzpicture}[>=Latex, scale=1.2]
            % Plano horizontal
            \draw[thick] (-2,0) -- (1,1) -- (3.7,1) -- (0.7,0) -- cycle;
            
            % Agujero
            \draw (0.5,0.5) circle (1.5pt);
            
            % Hilo y Partícula en el plano
            \draw[thick] (0.5,0.5) -- (1.8,0.7) node[midway, above] {$r$};
            \shadedraw[ball color=gris1!30] (1.8,0.7) circle (3pt) node[right, xshift=2pt] {$m$};
            
            % Resorte vertical
            \draw[decoration={aspect=0.3, segment length=3pt, amplitude=3pt, coil}, decorate] 
                (0.5,-0.3) -- (0.5,-1.2);
                \draw (0.5,-0.3)--(0.5,0.5);
            
            % Soporte resorte
            \fill[pattern=north east lines] (0.2,-1.4) rectangle (0.8,-1.2);
            \draw (0.2,-1.2) -- (0.8,-1.2);
            
            % Etiqueta constante k
            \node[right] at (0.6,-0.6) {$k$};
        \end{tikzpicture}
        \vfill
        \small (a) Perspectiva del sistema.
    \end{minipage}
    \hfill
    % --- VISTA B: PLANTA (SUPERIOR) ---
    \begin{minipage}{0.48\textwidth}
        \centering
        \begin{tikzpicture}[>=Latex, scale=1.2]
            % Ejes de referencia
            \draw[dashed, gris1!60] (0,0) -- (2,0);
            \draw[dashed, gris1!60] (0,0) -- (0,1.5);
            
            % Trayectoria y posición
            \draw[thick] (0,0) -- (30:1.8) node[midway, above left] {$r$};
            \shadedraw[ball color=gris1!30] (30:1.8) circle (3.5pt);
            
            % Ángulo theta
            \draw[->] (0.6,0) arc (0:30:0.6);
            \node at (0.8, 0.2) {$\theta$};
            
            % Vector velocidad inicial v0
            \draw[->, azul, thick] (30:1.8) -- +(120:0.8) node[above] {$\vec{v}_0$};
            
            % Origen (agujero)
            \fill (0,0) circle (1.5pt) node[below left] {$O$};
            
            % Gravedad (indicada saliendo/entrando al plano)
            \node[draw, circle, inner sep=1pt] at (2, 1.2) {$\times$};
            \node[right] at (2.1, 1.2) {$\mathbf{g}$};
        \end{tikzpicture}
        \vfill
        \small (b) Vista de planta y condiciones iniciales.
    \end{minipage}

    \caption{Sistema de partícula sobre plano horizontal con vínculo elástico vertical.}
    \label{fig:particula_resorte_vertical}
\end{figure}